\section{Proofs for the instance family}\label{app:family}

Throughout, $U\in\mathbb{R}^{n\times r}$ has rows $u_i$, $\Xs=UU^{\mathsf T}$,
$P_r$ is the orthogonal projector onto $\range U=\range\Xs$, and
$Q_0Q_0^{\mathsf T}=I-P_r$ (because $\Xs$ is symmetric, $\ker\Xs$ is the
orthogonal complement of $\range\Xs$).

\subsection{Proof of \cref{prop:det}}
With $\mu=c\,e$, we have $W=cQ_0Q_0^{\mathsf T}=c(I-P_r)$. For $i\ne j$, the
diagonal correction $\Diag(\diag W)$ in~\eqref{eq:construction} does not
affect the entry $(i,j)$, and $I_{ij}=0$, so
\[
  G_{ij}=\Xs_{ij}-W_{ij}=\Xs_{ij}+c\,(P_r)_{ij}.
\]
By the triangle inequality, $|G_{ij}|\le|\Xs_{ij}|+c|(P_r)_{ij}|\le a+cb$. If
$a<1$ and $c\le(1-a)/b$, then $a+cb\le1$. \qed

\begin{remark}
The bound is sufficient but not necessary. Cancellation between $\Xs_{ij}$ and
$c(P_r)_{ij}$, and the fact that $a$ and $b$ are generally attained at
different pairs $(i,j)$, both make it conservative. For this reason the
generator checks $|G_{ij}|\le1$ entry by entry.
\end{remark}

\subsection{Proof of \cref{prop:sat}}
Fix $r\ge2$ and $\epsilon\in(0,1)$. Because $S^{r-1}$ is compact, it can be
covered by finitely many caps
$C_1,\dots,C_{N}$, $N=N(\epsilon)$, each of angular radius $\epsilon/2$. Any
two points in the same cap make an angle of at most $\epsilon$, so their inner
product is at least $\cos\epsilon\ge1-\epsilon^2/2$. Once $n>N$, the
pigeonhole principle forces two of $u_1,\dots,u_n$ into a common cap, so
$a\ge1-\epsilon^2/2$ deterministically. Since $\epsilon$ is arbitrary and
$a\le1$ always, $a\to1$ along every realization. (A probabilistic argument is
needed only for a \emph{rate}; the limit itself follows from pigeonhole.)
\qed

\begin{remark}[What this does not give]
It is tempting to conclude that a fixed $c>0$ must eventually violate
\cref{prop:det}. That does not follow. At fixed $r$ the projector bound $b$
also tends to zero, so $a+cb\le1$ can persist: with $a_n=1-2/n$ and
$b_n=1/n$ one has $a_n\to1$ while $a_n+cb_n<1$ for every $n$ at $c=1$. A claim
in either direction needs the relative rates of $1-a$ and $b$, which we do not
establish here.
\end{remark}

\subsection{Proof of \cref{prop:hd}}
Fix $i\ne j$ and condition on $u_i$. Since $u_j$ is uniform on $S^{r-1}$ and
independent of $u_i$, the uniform distribution on the sphere is subgaussian:
by~\cite[Thm.~3.4.5]{vershynin2026high}, for every unit vector $v$ and every
$t\ge0$,
\[
  \Pr\bigl\{u_j^{\mathsf T}v\ge t\bigr\}\le 2\exp(-t^2r/2).
\]
Applying this with $v=u_i$, and again to $-u_j$, gives
$\Pr\{|u_i^{\mathsf T}u_j|\ge t\}\le4\exp(-t^2r/2)$; the bound is unchanged
after averaging over $u_i$. Taking a union bound over the fewer than $n^2/2$
unordered pairs,
\[
  \Pr\{a\ge t\}\;\le\;2n^2\exp(-t^2r/2).
\]
Setting the right-hand side equal to $\delta$ and solving for $t$ gives
$t=\sqrt{(2/r)(2\log n+\log(2/\delta))}$, which is the stated bound. If
$r/\log n\to\infty$ then $t\to0$ for any fixed $\delta$. \qed

\begin{remark}[On $b$]
We do not state a companion bound for $b=\max_{i\ne j}|(P_r)_{ij}|$. Writing
$P_r=VV^{\mathsf T}$ with $V$ having orthonormal columns, $(P_r)_{ij}$ is an
inner product of two rows of $V$, whose joint distribution is not that of two
independent uniform vectors; a bound would need a model for $\range U$ (for
instance a Haar-random subspace) together with a justification that the
row-normalized Gaussian $U$ induces it. Since the released family lies in the
fixed-rank regime of \cref{prop:sat}, where \cref{prop:hd} does not apply, and
since the generator screens $|G_{ij}|\le1$ entrywise in any case, we leave
this open.
\end{remark}

\begin{remark}[Scope]
These propositions concern $|G_{ij}|\le1$ only. Indefiniteness of $G$
($\lambda_{\min}(G)<0$) is not implied and is checked separately. The
released family ($r\in\{5,20,50\}$, $n=100$) is in the fixed-rank regime of
\cref{prop:sat}, not the high-dimensional regime of \cref{prop:hd}.
\end{remark}

\subsection{Mixed $\mu$}
For the released $\mu=(\mu_{\mathrm{bulk}}e_{k-m},\ \delta e_m)$, write
$Q_0=[Q_b\ Q_\delta]$. Then
$W=\mu_{\mathrm{bulk}}Q_bQ_b^{\mathsf T}+\delta Q_\delta Q_\delta^{\mathsf T}
=\mu_{\mathrm{bulk}}(I-P_r)-(\mu_{\mathrm{bulk}}-\delta)Q_\delta Q_\delta^{\mathsf T}$.
The off-diagonal entries therefore pick up an extra term
$(\mu_{\mathrm{bulk}}-\delta)(Q_\delta Q_\delta^{\mathsf T})_{ij}$, whose size
depends on the choice of basis for $\ker\Xs$ used to select $Q_\delta$. A
bound analogous to \cref{prop:det} follows with $b$ replaced by
$b+\max_{i\ne j}|(Q_\delta Q_\delta^{\mathsf T})_{ij}|$.

\begin{remark}[Open]
We do not characterize $\max_{i\ne j}|(Q_\delta Q_\delta^{\mathsf T})_{ij}|$.
It depends on which orthonormal basis of $\ker\Xs$ is chosen, and the
generator fixes that basis through the eigendecomposition it computes, so the
quantity is not determined by $(n,r,m,\delta)$ alone. A basis-free bound would
need either a canonical choice of $Q_\delta$ or a bound uniform over
admissible bases. We therefore state \cref{prop:det} for uniform $\mu$ only,
and the generator screens the mixed-$\mu$ instances entrywise rather than
relying on a bound. This is why the released family is screened rather than
certified.
\end{remark}
